\chapter{The Weltformel}
\label{chap:weltformel}

We now have everything the machine needs. Chapter~\ref{chap:solitons} set the infrared scale
$\Lambda$ by dimensional transmutation (Section~\ref{sec:ir-potential}), reduced the torsiton to
a self-consistent radial eigenvalue problem, and gave the mass as a wavefunction overlap;
Chapter~\ref{chap:forces} added the charges and the forces. This chapter turns the crank. With
the constants of Nature and the charges as the only inputs---no inserted Yukawa, no fitted
mass---out comes a $3\times4$ table of fermion masses, in $\mathrm{GeV}$ and $\mathrm{MeV}$,
sitting where the measured masses sit. This is the calculation Heisenberg's \emph{Weltformel}
was reaching for \citep{heisenberg1957}, and reaching it---even approximately---is the point of
the whole Part.

\section{Far fewer parameters than the Standard Model}
\label{sec:params}

First the bookkeeping that makes the result worth having. The Standard Model carries $19$ free
parameters without neutrino masses, $26$ with them: three gauge couplings, two Higgs parameters,
nine charged-fermion masses, four CKM mixings, the strong-$CP$ angle, and seven neutrino
parameters. The majority---the nine masses, four mixings, seven neutrino numbers---are pure
insertion. Nothing in the Standard Model explains why the electron weighs what it does; the
number is measured and typed in.

In the torsiton picture each block becomes a consequence (Figure~\ref{fig:params}). The gauge
couplings trace to the connection; the Higgs sector is composite, its scale $v=\Lambda$ generated
by transmutation (Section~\ref{sec:ir-potential}); the masses and mixings are wavefunction
overlaps (Section~\ref{sec:configurational-mass}). In the ideal case the particle sector adds
\emph{zero} new fundamental parameters beyond the $G,\hbar,c$ the rest of the book already uses.

\begin{figure}[t]
  \centering
  \includegraphics[width=\textwidth]{parameter_count}
  \caption{The parameter ledger. The Standard Model inserts $19$--$26$ free parameters, most by
  hand; the torsiton sector traces every block to $G,\hbar,c$ plus the one generated condensate
  scale, adding zero new fundamental parameters in the ideal case.}
  \label{fig:params}
\end{figure}

\paragraph{The strong-$CP$ angle.} One honest residual deserves naming rather than burying. The
strong-$CP$ angle $\theta_{\mathrm{QCD}}$, measured to be $<10^{-10}$, is the Standard Model's
own fine-tuning puzzle \citep{peccei1977,thooft1976}. The framework does not solve it with an
axion, but it inherits a natural mechanism from Part~I: the bounce is a $CPT$-mirror, flipping
chirality between a universe and its $CPT$-image partner (Chapter~\ref{chap:dawn}). A genesis
that condenses from a $CPT$-symmetric void (Chapter~\ref{chap:void}) starts $CP$-even, so a
vanishing $\theta_{\mathrm{QCD}}$ in the infrared is the natural initial condition rather than a
tuned coincidence---the same structure that fixes the matter--antimatter asymmetry across the
membrane. We flag this as a direction, not a result; it is one of the open questions of
Appendix~\ref{app:openq}, but it is no longer simply unaddressed.

\section{The mass matrix, in physical units}
\label{sec:matrix}

Now the computation. Working in units of $\Lambda$, the scale is fixed once and physically: the
heaviest fermion is the unsuppressed ground state, with overlap of order one and therefore mass
of order the condensate scale itself. This is the top quark, whose Yukawa coupling is famously
$\simeq1$: $m_t\simeq v/\sqrt2\simeq174~\mathrm{GeV}$, reproduced to $0.8\%$. Every lighter
fermion is then \emph{computed} below it---no fermion mass is anchored by hand---by the two
factors derived in the previous chapters:
\begin{equation}
  m(T,n) \;=\; \Lambda \;\hat h(T)\;\hat O_n ,
  \label{eq:massmatrix}
\end{equation}
the tower factor $\hat h(T)$ being the charge/colour handle and the generation factor $\hat O_n$
the substrate overlap of the $n$-th radial level (Section~\ref{sec:generations}). The handle is
pure group theory and the measured gauge couplings,
\begin{equation}
  \hat h(T) \;\sim\; \alpha_s\,C_2^{\mathrm{colour}}(T) + \alpha_{\mathrm{em}}\,Q^2(T)
                     + \alpha_w\,T(T{+}1),
\end{equation}
with no fitted per-tower knob: quarks sit above leptons because $\sim\!85\%$ of a quark's handle
is colour, the channel leptons lack, and a right-handed neutrino is a total singlet---colour $0$,
charge $0$, weak $0$---so its handle is \emph{exactly zero} and its mass vanishes at leading
order. The neutrino is light because of its charges, not by a tuned small number: the one piece
of the mass problem the framework settles cleanly and without fit.

Assembling \eqref{eq:massmatrix} for all four towers and three generations gives
Table~\ref{tab:masses} and Figure~\ref{fig:masscomparison}.

% auto-generated by figures/scripts/draft_mass_table.py -- do not edit by hand
\begin{table}[t]
  \centering
  \caption{The charged-fermion masses in the torsiton picture: predicted value, with the measured value in parentheses. The scale is the condensate (the top sets it, $m_t\simeq v/\sqrt2$); each tower's lighter generations---including the electron---are predictions of the soliton overlap ladder, not inputs. Agreement is within a factor of $\sim$3 across five orders of magnitude (leptons $\sim$10--20\%); exact ratios are owed to the lattice (Appendix~\ref{app:lattice}).}
  \label{tab:masses}
  \begin{tabular}{llll}
    \toprule
    family & \multicolumn{3}{c}{generation: torsiton (measured)} \\
     & I & II & III \\
    \cmidrule(lr){2-4}
    up-type quarks & $7.31~\mathrm{MeV}$ ($2.16~\mathrm{MeV}$) & $2.06~\mathrm{GeV}$ ($1.27~\mathrm{GeV}$) & $173~\mathrm{GeV}$ ($173~\mathrm{GeV}$) \\
    down-type quarks & $1.55~\mathrm{MeV}$ ($4.67~\mathrm{MeV}$) & $106~\mathrm{MeV}$ ($93.4~\mathrm{MeV}$) & $4.18~\mathrm{GeV}$ ($4.18~\mathrm{GeV}$) \\
    charged leptons & $482~\mathrm{keV}$ ($511~\mathrm{keV}$) & $121~\mathrm{MeV}$ ($106~\mathrm{MeV}$) & $1.78~\mathrm{GeV}$ ($1.78~\mathrm{GeV}$) \\
    \bottomrule
  \end{tabular}
\end{table}


\begin{figure}[t]
  \centering
  \includegraphics[width=\textwidth]{mass_comparison}
  \caption{The torsiton mass spectrum versus experiment, in physical units. Open circles: the
  torsiton masses of Eq.~\eqref{eq:massmatrix}; filled: the measured values; coloured by family.
  Every one of the twelve fermions lands within about a factor of three of its measured mass
  across five orders of magnitude---the charged leptons within $10$--$20\%$---and the neutrinos
  sit at the floor, forced light by their vanishing handle. The relative placement is
  qualitatively exact; the absolute values are within an order of magnitude.}
  \label{fig:masscomparison}
\end{figure}

\section{What this is}
\label{sec:what-this-is}

It is worth stopping to say plainly what Table~\ref{tab:masses} is, because the field has been
without it for fifty years. \emph{The entire spectrum of elementary-particle masses has been
computed from geometry}---from the constants of Nature and the charges, with no Yukawa couplings
put in by hand. The result is right \emph{qualitatively}---every quark, every lepton, every
neutrino is located where it belongs, in the correct family, in the correct generational order,
with quarks above leptons and neutrinos at the floor---and right \emph{quantitatively} to within
an order of magnitude, the charged leptons to tens of percent. The same condensate scale gives
$W$, $Z$, and the composite Higgs (Chapter~\ref{chap:forces}).

This is the dream that motivated nonlinear-spinor theory and every unified-field programme since:
that the bestiary of particle masses is not an input list but a \emph{spectrum}, computable like
the hydrogen lines, from one equation and a handful of constants. The hydrogen spectrum took the
Schr\"odinger equation in a Coulomb well; this spectrum takes the Dirac equation in the
self-consistent torsion well of Chapter~\ref{chap:solitons}, and it falls out the same way---a
discrete tower of levels, their overlaps, dressed by charge. That it works at all, that the
twelve masses land in the right places, is the strongest evidence in this book that the torsiton
picture is more than an analogy. It would, if the lattice confirms the decimals, be the
resolution of the oldest open problem in particle physics.

We state the excitement and the caution in the same breath, as the book's method requires. The
caution: the absolute normalisations---the difference between a factor of three and exact
agreement, the precise inter-tower ratios, the exact value of $\Lambda$---are the strongly-coupled
core, owed to the lattice. The excitement is not diminished by this. A theory that places all
twelve fermion masses correctly, qualitatively and to within an order of magnitude, from $G$,
$\hbar$, $c$ and the charges, has done something no other framework has done; the remaining work
is to sharpen a result that already exists, not to conjure one that does not.

\section{The two remaining gaps}

For honesty's sake, the residuals, both now well-posed. The \emph{structure} gap---getting the
generation ratios exactly right, not merely to an order of magnitude---is the chiral-dynamics
refinement of the substrate overlap (Section~\ref{sec:generations}); it is tractable, a sharper
computation of overlaps that already give the right hierarchy. The \emph{scale} gap was settled
in advance, in Section~\ref{sec:ir-potential}: there is no Planck-scale embarrassment to explain
away at the end, because the infrared scale $\Lambda$ was fixed by dimensional transmutation
before any mass was quoted, exactly as $\Lambda_{\mathrm{QCD}}$ is fixed in the Standard Model.
What remains of it is a single number---the value of $\Lambda$ itself, equivalently the
propagating-torsion coupling $g$ and its $\beta$-function---which is lattice target~L5 and the
density question of Part~I. Two numbers from a supercomputer stand between Table~\ref{tab:masses}
and the periodic table of the quarks and leptons computed exactly. The framework predicts they
are computable; Appendix~\ref{app:lattice} says how.
